\section{Shape conditions and fitted parameters}\label{sec:transport}

We use the fitted functions $\lambda(h,r),\kappa(h,r)$ defined in
Equation~\eqref{eq:fit-identities} for a smooth profile $\phi$,
and retain $k=\phi'$, $D=\partial_h-\alpha\partial_r$ and the
model derivatives evaluated at the fit. We will prove two comparisons:
the sign of $\phi'''$ controls the width derivative of fitted curvature,
while the sign of $((\phi'-d)/u)'$, for a fixed real $d$, controls
the fitted slope intercept $\lambda-h\kappa$.
These are the shape conditions used in Section~\ref{sec:application}.
The transport law is already available from
Lemma~\ref{lem:fit-transport}; its evolution starts from the
small-width data established next.

\begin{lemma}[The small-width fit]\label{lem:fit-small-width}
At each positive height, locally uniformly in that height as $r$ tends
to zero,
\begin{equation}\label{eq:fit-small-width}
 \begin{aligned}
 \kappa(h,r)&=\phi''(h)+\frac{2}{7}\phi'''(h)r^2+O(r^3),\\
 \kappa_r(h,r)&=\frac{4}{7}\phi'''(h)r+O(r^2),\\
 \lambda(h,r)&=k(h)-\frac{1}{35}\phi'''(h)r^4+O(r^5).
 \end{aligned}
\end{equation}
More precisely, for every $h_0>0$ there are smooth functions $A$ and
$B$ on a neighborhood of $(h_0,0)$ for which, at positive $r$ in that
neighborhood,
\begin{equation}\label{eq:fit-smooth-remainders}
 \begin{aligned}
 \lambda(h,r)&=k(h)-\frac{1}{35}\phi'''(h)r^4+r^5 A(h,r),\\
 \kappa(h,r)&=\phi''(h)+\frac{2}{7}\phi'''(h)r^2+r^3 B(h,r).
 \end{aligned}
\end{equation}
Thus the derivative estimate in Equation~\eqref{eq:fit-small-width}
comes from a smooth remainder.
\end{lemma}

\begin{proof}
Fix $h_0>0$. Use the smooth width-one profile family
$g(h,\rho,s)$ of Equation~\eqref{eq:arc-scaling-profile}, with signed
auxiliary parameter $\rho$ near zero and $h$ near $h_0$.
Lemma~\ref{lem:arc-scaling} provides smooth endpoint data at scaled
width one for this family and for the finite-parameter perturbations
below. In particular, the local implicit inversion that fixes the
width is performed near the same transverse parabolic arc.

For a scaled profile $\psi$, let $\mathcal A(\psi)$ denote the pair
consisting of its midpoint and its midpoint derivative with respect
to scaled width, evaluated at base height $0$ and width $1$.  Only the
smooth finite-parameter families of profiles just described are used
in this notation.  For the quadratic profile
$\psi(s)=a_1s+a_2s^2/2$, the derivative of this pair at $(a_1,a_2)=(0,0)$
is, by Equation~\eqref{eq:parabolic-variations},
\[
 J=\begin{pmatrix}
       1/3&1/15\\
       2/3&4/15
    \end{pmatrix}.
\]
It is nonsingular.  The cubic profile variation $s^3/6$ gives the
column
\[
 b=\binom{1/105}{6/105},\qquad
 J^{-1}b=\binom{-1/35}{2/7}.
\]

Taylor expansion of Equation~\eqref{eq:arc-scaling-profile}, with a
smooth remainder on the fixed scaled-height interval, gives
\begin{equation}\label{eq:fit-scaled-taylor}
 g(h,\rho,s)
 =\rho k(h)s+\rho^3\phi''(h)\frac{s^2}{2}
   +\rho^5\phi'''(h)\frac{s^3}{6}
   +\rho^7\mathcal R(h,\rho,s),
\end{equation}
where $\mathcal R$ is smooth.  Write
\[
 g_2(h,\rho,s)=\rho k(h)s+\rho^3\phi''(h)\frac{s^2}{2},
 \qquad
 \eta(h,\rho,s)=\phi'''(h)\frac{s^3}{6}
                    +\rho^2\mathcal R(h,\rho,s).
\]
Then $g=g_2+\rho^5\eta$.  Apply the local inverse defined by $J$ to
the data of the smooth family $g_2+\varepsilon\eta$, with
$(h,\rho,\varepsilon)$ near $(h_0,0,0)$.  Denote its fitted scaled
coefficients by $a_1(h,\rho,\varepsilon)$ and
$a_2(h,\rho,\varepsilon)$.  At $\varepsilon=0$ these are exactly
$\rho k(h)$ and $\rho^3\phi''(h)$, respectively.  The integral form of
Taylor's formula in $\varepsilon$ shows that
\[
 \binom{a_1(h,\rho,\rho^5)-\rho k(h)}
       {a_2(h,\rho,\rho^5)-\rho^3\phi''(h)}
   =\rho^5 q(h,\rho)
\]
for a smooth vector-valued function $q$.  Differentiating at
$\rho=\varepsilon=0$ in the profile-variation direction yields
\[
 q(h,0)=\phi'''(h)J^{-1}b
       =\phi'''(h)\binom{-1/35}{2/7}.
\]
Another smooth Taylor factorization in $\rho$ therefore writes
\[
 q(h,\rho)=\phi'''(h)\binom{-1/35}{2/7}
                    +\rho\binom{A(h,\rho)}{B(h,\rho)}.
\]

For $\rho=r>0$, the scaled actual data are $(M(h,r)/r,v(h,r))$.
Scaling the quadratic fit in Equation~\eqref{eq:fit-identities} gives
scaled coefficients $(r\lambda,r^3\kappa)$.  Uniqueness in
Theorem~\ref{thm:quadratic-fit} identifies them with the local
coefficients just constructed.  Dividing the two corrections by $r$
and $r^3$, respectively, proves
Equation~\eqref{eq:fit-smooth-remainders}.  Differentiating its second
identity at fixed $h$ gives
\[
 \kappa_r(h,r)=\frac47\phi'''(h)r
                 +3r^2B(h,r)+r^3B_r(h,r).
\]
This proves all the locally uniform estimates in
Equation~\eqref{eq:fit-small-width}.
\end{proof}

\begin{proposition}[A strict third derivative]\label{prop:third-derivative}
If $\phi'''(u)>0$ for every $u>0$, then, for every $h,r>0$,
\[
 \lambda(h,r)<k(h),\qquad \kappa_r(h,r)>0.
\]
If $\phi'''(u)<0$ for every $u>0$, both inequalities reverse:
\[
 \lambda(h,r)>k(h),\qquad \kappa_r(h,r)<0.
\]
\end{proposition}

\begin{proof}
Assume first that $\phi'''>0$.  Fix a fitted arc with base height $h$
and width $r$, and let its peak height be $t$.  Parametrize its
characteristic by outward width $s\in(0,r]$, writing the base height
as $h(s)$, with $h(r)=h$ and $h(s)\to t$ as $s\downarrow0$.  Put
\[
 \theta(s)=\operatorname{arctanh}\frac{M(h(s),s)}s,
 \qquad
 L(s)=\operatorname{arctanh}v(h(s),s),
 \qquad c(s)=\phi''(h(s)).
\]
The characteristic equations in
Equation~\eqref{eq:arc-characteristics} give
\begin{equation}\label{eq:actual-curvature-decrease}
 c'(s)=-s\operatorname{sech}^2\theta(s)\,
                       \phi'''(h(s))<0\qquad(s>0).
\end{equation}
In particular $c$ extends continuously to $s=0$ with $c(0)=\phi''(t)$
and is strictly decreasing on $[0,r]$.

In the following comparison, hold the curvature equal to the final
fitted value $\kappa=\kappa(h,r)$.  Let $J_\kappa$ and
$\Theta_\kappa(s,z)$ be the slope field and the reference angle
family of Section~\ref{sec:fitting}.  Define
\[
 E(s)=k(h(s))-J_\kappa(s,\theta(s)),\qquad
 a(s)=(J_\kappa)_\theta(s,\theta(s))>0,
\]
and
\[
 f(s)=-s\operatorname{sech}^2\theta(s)\,(c(s)-\kappa).
\]
Lemma~\ref{lem:curvature-comparison}, specifically
Equation~\eqref{eq:curvature-error}, gives
\begin{equation}\label{eq:third-derivative-error}
 E'(s)+a(s)E(s)=f(s),\qquad \lim_{s\downarrow0}E(s)=0.
\end{equation}
The integral formula in Equation~\eqref{eq:fit-error-integral}
applies with this forcing $f$. Its positive kernel preserves the
sign of $f$, even though $a$ may be singular at the peak.

The strict decrease of $c$ gives the following exhaustive sign cases.
If $\kappa\ge c(0)$, then $f(s)>0$ for every $s>0$, and
Equation~\eqref{eq:fit-error-integral} gives $E(s)>0$.
If $\kappa\le c(r)$, then $f(s)<0$ for $0<s<r$, and $E(s)<0$ for
$0<s\le r$.  In the remaining case there is a unique
$s_0\in(0,r)$ with $c(s_0)=\kappa$.  Here $f<0$ before $s_0$ and
$f>0$ after it.  The same integral formula gives $E<0$ on $(0,s_0]$,
whereas
\[
 \frac{d}{ds}\left[
 E(s)\exp\left(\int_{s_0}^s a(q)\,dq\right)\right]
 =f(s)\exp\left(\int_{s_0}^s a(q)\,dq\right)>0
 \qquad(s>s_0).
\]
Consequently $E$ either remains negative up to the final width,
possibly having $E(r)=0$, or crosses zero exactly once before $r$
and is strictly positive thereafter.  It cannot vanish on a
nontrivial interval: Equation~\eqref{eq:third-derivative-error} would
then force $c=\kappa$ on that interval, contrary to
Equation~\eqref{eq:actual-curvature-decrease}.

Let $z(s)$ be the reference label determined by
\[
 \theta(s)=\Theta_\kappa(s,z(s)).
\]
By Equation~\eqref{eq:label-derivative},
\[
 z'(s)=\frac{E(s)}{\partial_z\Theta_\kappa(s,z(s))},
 \qquad \partial_z\Theta_\kappa(s,z)>0.
\]
The constant label $z_* = z(r)$ describes the quadratic reference
whose final angle agrees with the actual one.  Its final slope is
$\lambda(h,r)$, by the construction of the fit.  The fit also makes
the final values of $v$, and hence of $L$, equal.  Therefore
Equation~\eqref{eq:arc-characteristics} gives
\begin{equation}\label{eq:matched-angle-integrals}
 \int_0^r\frac{\sinh(2\theta(s))}{s}\,ds
 =\int_0^r\frac{\sinh(2\Theta_\kappa(s,z_*))}{s}\,ds.
\end{equation}
Both integrals converge: the two angles are $O(s)$ near their smooth
peaks.

If $z(s)$ were monotone and nonconstant on $(0,r]$, monotonicity of
$\Theta_\kappa$ in its label would order the two integrands in
Equation~\eqref{eq:matched-angle-integrals} in one direction, with
strict inequality on a nonempty open subinterval.  This contradicts
their equal integrals.  The first two sign cases above therefore
cannot occur.  In the two-sign forcing case, $E$ must cross zero
strictly before $r$: otherwise $z$ is strictly decreasing on $(0,r)$,
even if $E(r)=0$, giving the same contradiction.  A constant label is
also impossible, since it would give $E\equiv0$ and hence constant $c$.
It follows that
\begin{equation}\label{eq:fitted-slope-strict}
 E(r)=k(h)-\lambda(h,r)>0.
\end{equation}

It remains to propagate the sign of the width derivative of fitted
curvature.  As a function of independent model parameters, set
\[
 C(\lambda,\kappa,r)=\frac{R(\lambda,\kappa,r)}
                          {G(\lambda,\kappa,r)}.
\]
Theorem~\ref{thm:model-inequality} asserts that its explicit partial
derivative satisfies $C_r\le0$ when $\lambda$ and $\kappa$ are held
fixed.  Differentiating the second identity in
Equation~\eqref{eq:fit-transport} at fixed $h$, we must also
account for
\[
 \partial_r(D\kappa)=D(\kappa_r)-\alpha_r\kappa_r.
\]
Using $\lambda_r=-Q\kappa_r/P$ gives
\begin{equation}\label{eq:curvature-derivative-transport}
 D(\kappa_r)=\mathcal T(h,r)\kappa_r
                 -C_r(\lambda,\kappa,r)(\lambda-k(h)),
\end{equation}
where
\begin{equation}\label{eq:curvature-derivative-coefficient}
 \mathcal T
  =\alpha_r+\frac{CQ}{P}
     -\left(C_\kappa-\frac{Q}{P}C_\lambda\right)(\lambda-k(h)).
\end{equation}
Here $\alpha_r$ is the derivative of the actual function
$\alpha(h,r)$, whereas $C_\lambda,C_\kappa,C_r$ are model partial
derivatives evaluated at the fit.  All terms in $\mathcal T$ are
smooth at positive height and width.  The dependence of $\alpha$ on
$M$ is thus included in the coefficient of $\kappa_r$; it creates no
additional independent source.

On the fixed characteristic, write
$w(s)=\kappa_r(h(s),s)$.  Equation~\eqref{eq:fit-small-width},
locally uniformly near its peak $t$, gives
\[
 w(s)=\frac47\phi'''(h(s))s+O(s^2)>0
\]
for all sufficiently small positive $s$.  Equations
\eqref{eq:fitted-slope-strict} and
\eqref{eq:curvature-derivative-transport} give, along the
characteristic,
\[
 w'(s)=-\frac{\mathcal T}{\alpha}w(s)
              +\frac{C_r}{\alpha}(\lambda-k),
 \qquad \frac{C_r}{\alpha}(\lambda-k)\ge0.
\]
Choose a sufficiently small positive $s_1$ with $w(s_1)>0$.  On
$[s_1,r]$ all coefficients are continuous.  Multiplication by the
positive integrating factor
\[
 \mu(s)=\exp\left(\int_{s_1}^s\frac{\mathcal T}{\alpha}\,dq\right)
\]
gives $(\mu w)'\ge0$.  Thus $w(r)>0$, proving $\kappa_r(h,r)>0$.
Since the chosen fit was arbitrary, both assertions hold for every
$h,r>0$.

Finally, reflection of the endpoint coordinate changes $\phi$ to
$-\phi$ and changes $(M,v)$ to $(-M,-v)$.  Uniqueness of the quadratic
fit changes $(\lambda,\kappa)$ to $(-\lambda,-\kappa)$.  Applying the
proved case to $-\phi$ proves the two reversed inequalities when
$\phi'''<0$.
\end{proof}

\begin{proposition}[The fitted slope intercept]\label{prop:intercept}
Let $d\in\mathbb R$.  If
\[
 \frac{d}{du}\left(\frac{k(u)-d}{u}\right)>0
 \qquad\text{for every }u>0,
\]
then
\[
 \lambda(h,r)-h\kappa(h,r)<d
 \qquad\text{for every }h,r>0.
\]
If the displayed derivative is strictly negative everywhere, then
$\lambda(h,r)-h\kappa(h,r)>d$ everywhere.
\end{proposition}

\begin{proof}
Assume the displayed derivative is positive and set
\[
 I(h,r)=\lambda(h,r)-h\kappa(h,r)-d.
\]
Fix a characteristic with peak height $t$.  By
Equation~\eqref{eq:fit-small-width}, its limit at width zero is
\[
 \lim_{s\downarrow0}I(h(s),s)
 =k(t)-t\phi''(t)-d
 =-t^2\left.\frac{d}{du}\left(\frac{k(u)-d}{u}\right)\right|_{u=t}
 <0.
\]
Thus $I$ is strictly negative near the peak.

Consider any positive-height point on the characteristic at which
$I=0$.  At that point $\lambda=d+h\kappa$.  We claim that
\begin{equation}\label{eq:intercept-zero-curvature}
 \kappa>\frac{k(h)-d}{h}.
\end{equation}
Indeed, if the reverse weak inequality held, then for every $u>h$
the hypothesis would give
\[
 \frac{k(u)-d}{u}>\frac{k(h)-d}{h}\ge\kappa,
 \qquad
 k(u)>d+u\kappa=\lambda+(u-h)\kappa.
\]
Integrating from $h$ to $u$ shows
\[
 \phi(u)-\phi(h)
   >\lambda(u-h)+\frac{\kappa}{2}(u-h)^2
 \qquad(u>h).
\]
The strict translated-profile comparison of Lemma~\ref{lem:variation}
then makes the actual midpoint at the fixed base height and width
strictly larger than the quadratic midpoint.  This contradicts
Equation~\eqref{eq:fit-identities}, proving
Equation~\eqref{eq:intercept-zero-curvature}.  In particular
$\lambda=d+h\kappa>k(h)$ at any zero of $I$.

Equation~\eqref{eq:fit-transport} gives the exact identity
\begin{equation}\label{eq:intercept-transport}
 DI=\frac{S+hR}{G}\bigl(\lambda-k(h)\bigr).
\end{equation}
At a zero of $I$ its right-hand side is strictly positive.  The
outward characteristic derivative is therefore $-DI/\alpha<0$.
If $I$ had a first zero after its initially negative interval, its
outward derivative there would instead be nonnegative.  This
contradiction proves $I<0$ everywhere along the characteristic.
Every fit belongs to such a characteristic, so the assertion holds
for all $h,r>0$.

For a strictly negative derivative, apply the proved assertion to
$-\phi$ and $-d$, and use the reflection of fitted parameters from
the proof of Proposition~\ref{prop:third-derivative}.
\end{proof}

\begin{lemma}[Range of the fitted curvature]\label{lem:fitted-curvature-range}
Let an arc have positive base height $h$, peak height $t>h$, and
width $r$.  Its fitted curvature satisfies
\begin{equation}\label{eq:fitted-curvature-range}
 \min_{h\le u\le t}\phi''(u)
       \le\kappa(h,r)\le
 \max_{h\le u\le t}\phi''(u).
\end{equation}
If $\phi''$ is nonconstant on $[h,t]$, both inequalities are strict.
\end{lemma}

\begin{proof}
Write $m=\min_{[h,t]}\phi''$ and $b=\max_{[h,t]}\phi''$.
If the fitted $\kappa$ were less than $m$, the actual curvature would
be strictly larger than this constant reference curvature along the
whole characteristic.  Lemma~\ref{lem:curvature-comparison} would
then make the actual final $L$ strictly larger than the reference $L$
with the same final angle, contradicting the fit.  A fitted curvature
greater than $b$ gives the reversed contradiction.  This proves
Equation~\eqref{eq:fitted-curvature-range}.

If $\phi''$ is nonconstant and $\kappa=m$, the forcing in
Equation~\eqref{eq:curvature-error} is nonpositive and strictly
negative on a positive interior width interval.  Its integral
representation gives $E\le0$ everywhere and $E<0$ on a nonempty
interval.  Equation~\eqref{eq:label-derivative} makes the reference
label nonincreasing and nonconstant.  Relative to its final constant
label, the actual earlier angles are therefore no smaller and are
strictly larger on a nonempty open interval.  As in
Equation~\eqref{eq:matched-angle-integrals}, their $L$ integral is
strictly larger, a contradiction.  The case $\kappa=b$ reverses all
these inequalities.  Thus both bounds are strict whenever the
actual curvature is nonconstant.
\end{proof}
